\section{Local maximizers for the Faber--Krahn problem in the Gaussian STFT setting}

The purpose of this section is to establish a local rigidity statement for the
Faber--Krahn problem associated with the Gaussian STFT. The global problem,
solved in \cite{Nicola-Tilli}, asserts that, among measurable subsets of \(\C\)
of fixed measure, the centered disk uniquely maximizes the first eigenvalue of
the Gaussian STFT localization operator. Here we show that this rigidity is
already visible at the \emph{local} level: any set whose first eigenvalue cannot
be increased by small symmetric-difference perturbations must, up to a
translation, be a disk. This complements the quantitative stability result of
\cite{Gomez-Guerra-Ramos-Tilli} and provides the variational counterpart of the
inverse-spectral analysis carried out in the previous sections.

We work throughout in the Bargmann--Fock representation; we refer the reader to
\cite{Bargmann,Folland,Grochenig} for the relevant background. Thus, if
\(F\in \mathcal F^2(\C)\) is the Bargmann transform of \(f\in L^2(\R)\), then
\[
|\mathcal V_\varphi f(z)|^2 = |F(z)|^2 e^{-\pi |z|^2},
\]
up to the normalization already fixed in the previous sections. Accordingly, for
\(F\in \mathcal F^2(\C)\) we set
\[
u_F(z):=|F(z)|^2 e^{-\pi |z|^2}.
\]

For a measurable set \(\Omega\subset \C\), let
\[
\lambda_{1,\varphi}(\Omega)
:=
\sup_{\|f\|_2=1}\int_\Omega |\mathcal V_\varphi f(z)|^2\,dz.
\]
We say that \(\Omega\) is a \emph{local maximizer} for \(\lambda_{1,\varphi}\) under the
measure constraint if there exists \(\delta_0>0\) such that, whenever
\(E\subset \C\) is measurable, \(|E|=|\Omega|\), and
\[
\int_\C |\mathbf 1_\Omega-\mathbf 1_E|\,e^{-\pi |z|^2}\,dz<\delta_0,
\]
one has
\[
\lambda_{1,\varphi}(\Omega)\ge \lambda_{1,\varphi}(E).
\]

We now prove the local rigidity statement announced in the introduction. The
proof proceeds in six steps. Steps 1--3 are essentially the analogues, in our
setting, of standard arguments in the Faber--Krahn theory: \(\Sigma\) is forced
to be a level set of \(u_F\), the corresponding eigenvalue is simple, and a
first-variation identity ties \(F\) to a divergence-free vector field built from
holomorphic data. The crucial new ingredient lies in Step 4, where we show that
the \emph{derivative} \(F'\) is again a first eigenfunction. Combined with the
simplicity from Step 2, this forces \(F\) to satisfy a first-order linear ODE,
and hence to be a Gaussian. Once this is in place, Steps 5--6 use the Weyl
invariance of the Bargmann--Fock setting and a one-dimensional rearrangement
argument to identify \(\Sigma\) with a centered disk.

\begin{proof}[Proof of the local-maximizer theorem]
Let \(\Sigma\subset \C\) be a local maximizer for \(\lambda_{1,\varphi}\) under
the measure constraint, and let \(F\in \mathcal F^2(\C)\) be a normalized first
eigenfunction for the localization operator associated with \(\Sigma\).
Equivalently,
\begin{equation}\label{eq:eigenvalue-equation}
\int_\Sigma H(z)\overline{F(z)}\,d\mu(z)
=
\lambda_{1,\varphi}(\Sigma)\int_\C H(z)\overline{F(z)}\,d\mu(z)
\end{equation}
for every \(H\in \mathcal F^2(\C)\), where
\[
d\mu(z):=e^{-\pi |z|^2}\,dA(z).
\]

\medskip
\noindent
{\bf Step 1: \(\Sigma\) is a superlevel set of \(u_F\).}

We claim that there exists \(t_0\ge 0\) such that, up to sets of measure zero,
\begin{equation}\label{eq:superlevel-claim}
\Sigma=\{u_F>t_0\}.
\end{equation}

The strategy is a bathtub-style perturbation argument: if \(\Sigma\) were not a
superlevel set of \(u_F\), one could move a small piece of \(\Sigma\) where
\(u_F\) is small to a region where \(u_F\) is larger, increasing the
localization integral while preserving Lebesgue measure and remaining within
the local-maximizer neighborhood of \(\Sigma\).

Recall that \(F\) is admissible in the variational problem defining
\(\lambda_{1,\varphi}(\Sigma')\) for any measurable set \(\Sigma'\) of finite
measure. In particular,
\[
\lambda_{1,\varphi}(\Sigma')\ge \int_{\Sigma'} u_F\,dA,
\qquad
\lambda_{1,\varphi}(\Sigma)= \int_\Sigma u_F\,dA,
\]
the latter equality holding because \(F\) is the actual eigenfunction for
\(\Sigma\). It therefore suffices to find a competitor \(\Sigma'\) with
\(|\Sigma'|=|\Sigma|\), close to \(\Sigma\) in the sense of weighted symmetric
difference, and with \(\int_{\Sigma'} u_F\,dA>\int_\Sigma u_F\,dA\); that would
violate local maximality.

Suppose, for the sake of contradiction, that \eqref{eq:superlevel-claim} fails.
Choose, by the layer-cake representation, the unique \(t_0\ge 0\) such that
\[
A:=\{u_F>t_0\}\quad\text{satisfies}\quad |A|=|\Sigma|;
\]
then \(A\) and \(\Sigma\) do not coincide a.e., so both \(A\setminus\Sigma\) and
\(\Sigma\setminus A\) have positive Lebesgue measure (and in fact positive
weighted measure, since \(e^{-\pi|z|^2}>0\) everywhere).

We now perform the bathtub construction explicitly. By inner regularity of the
Lebesgue measure, choose an open subset
\(S'\subset A\setminus\Sigma\) and an open subset
\(T'\subset \Sigma\setminus A\) with \(|S'|,|T'|>0\). Replacing each by a
sufficiently small open subset, we may further assume
\(\int_{S'}d\mu,\,\int_{T'}d\mu<\delta_0/4\). Since the Lebesgue measure
restricted to a fixed open set has the intermediate-value property, we may
choose subsets \(S\subset S'\) and \(T\subset T'\) with the common Lebesgue
measure
\[
|S|=|T|=\tfrac12\min(|S'|,|T'|)>0.
\]
Then
\[
\int_\C (\mathbf 1_S+\mathbf 1_T)\,e^{-\pi |z|^2}\,dz<\tfrac{\delta_0}{2}<\delta_0.
\]
Define
\[
\Sigma':=(\Sigma\setminus T)\cup S.
\]
Then \(|\Sigma'|=|\Sigma|\) and
\[
\int_\C |\mathbf 1_\Sigma-\mathbf 1_{\Sigma'}|e^{-\pi |z|^2}\,dz
=\int_\C(\mathbf 1_S+\mathbf 1_T)e^{-\pi|z|^2}\,dz<\delta_0.
\]
On the other hand, by construction,
\[
u_F(z)>t_0 \quad \text{on } S\subset A,
\qquad
u_F(z)\le t_0 \quad \text{on } T\subset \Sigma\setminus A.
\]
Moreover, since \(u_F\) is continuous and \(S\subset A=\{u_F>t_0\}\) is open
with positive measure, we may by shrinking \(S\) further assume that there
exists \(\eta>0\) with \(u_F\ge t_0+\eta\) on \(S\). It then follows that
\[
\int_S u_F\,dA \ge (t_0+\eta)\,|S|,
\qquad
\int_T u_F\,dA \le t_0\,|T|=t_0\,|S|,
\]
and hence
\[
\int_{\Sigma'} u_F\,dA
=
\int_\Sigma u_F\,dA + \int_S u_F\,dA-\int_T u_F\,dA
\ge \int_\Sigma u_F\,dA + \eta\,|S|
>
\int_\Sigma u_F\,dA.
\]
Since \(F\) is admissible for the variational problem defining
\(\lambda_{1,\varphi}(\Sigma')\),
\[
\lambda_{1,\varphi}(\Sigma')
\ge \int_{\Sigma'} u_F\,dA
>
\int_\Sigma u_F\,dA
=
\lambda_{1,\varphi}(\Sigma),
\]
contradicting the local maximality of \(\Sigma\). This proves
\eqref{eq:superlevel-claim}.

In particular, since \(u_F\) is real-analytic outside the (discrete) zero set of
\(F\), \(u_F\) is constant on \(\partial \Sigma\) in the sense that there exists
\(c_\Sigma\ge 0\) with
\begin{equation}\label{eq:boundary-level-set-new}
u_F(z)=c_\Sigma
\qquad \text{for } z\in \partial \Sigma.
\end{equation}
We will see in Step 2 that \(F\) does not vanish on \(\overline\Sigma\), and in
particular \(c_\Sigma>0\).

\medskip
\noindent
{\bf Step 2: simplicity of the first eigenvalue.}

We claim that the first eigenvalue \(\lambda_{1,\varphi}(\Sigma)\) is simple.
The argument uses a modular-deformation principle: any two first eigenfunctions
have the same modulus on \(\partial\Sigma\) up to a constant, and a holomorphic
function with constant modulus on a connected open set must itself be constant.
This kind of argument plays a role also in the stability analysis of
\cite{Gomez-Guerra-Ramos-Tilli}.

Suppose, for contradiction, that \(G\in \mathcal F^2(\C)\) is another first
eigenfunction, linearly independent of \(F\). Step 1 applied to \(G\) shows that
\(\Sigma\) is also a superlevel set of \(u_G\); hence there exists
\(c'_\Sigma\ge 0\) such that
\[
u_G(z)=c'_\Sigma
\qquad\text{for } z\in \partial \Sigma.
\]
Equivalently,
\[
|F(z)|^2 e^{-\pi |z|^2}=c_\Sigma,
\qquad
|G(z)|^2 e^{-\pi |z|^2}=c'_\Sigma,
\qquad z\in \partial \Sigma.
\]

We next observe that first eigenfunctions are zero-free in the interior of
\(\Sigma\). Indeed, by Step 1 we have \(u_F>c_\Sigma\) on the (open) interior of
\(\Sigma\). Since
\[
u_F=|F|^2 e^{-\pi |z|^2}
\]
and the Gaussian factor \(e^{-\pi|z|^2}\) is strictly positive everywhere, this
yields \(|F|>0\) throughout the interior of \(\Sigma\); the same applies to
\(|G|\). In particular, \(c_\Sigma>0\) and \(c'_\Sigma>0\). Moreover, since
\(F,G\) are entire and continuous, \(|F|>0\) on a neighborhood of any compact
subset of \(\Sigma^\circ\); by Step 1, we may take a slight enlargement of
\(\Sigma\) (still a superlevel set of \(u_F\)) on which \(|F|,|G|>0\), so we may
work on a connected open set \(U\Supset \Sigma\) where both \(F\) and \(G\) are
zero-free.

The ratio \(F/G\) is holomorphic and nowhere zero on \(U\), and the function
\[
\log \left|\frac{F}{G}\right|
=
\frac12 \log \frac{|F|^2}{|G|^2}
=
\frac12 \log \frac{u_F}{u_G}
\]
is harmonic on \(U\): it is the real part of any local branch of
\(\log(F/G)\), which is well defined on simply connected subsets, and harmonic
in any case as the logarithm of the modulus of a zero-free holomorphic
function. Note that the Gaussian factors cancel inside the logarithm, so on
\(\partial \Sigma\) by \eqref{eq:boundary-level-set-new} this function equals
the constant \(\frac12\log(c_\Sigma/c'_\Sigma)\). By the maximum principle
applied separately on each connected component of \(\Sigma\), the function
\(\log|F/G|\) is constant on each component. Hence
\[
\left|\frac{F}{G}\right| \equiv \text{constant}
\qquad \text{on each connected component of } \Sigma.
\]
A nonconstant holomorphic function is open, so a holomorphic function with
constant modulus on a connected open set must itself be constant. Therefore
\[
F=\gamma_j G
\qquad \text{on each connected component } \Sigma_j\text{ of } \Sigma.
\]
The relation \(F=\gamma_j G\) is an identity between two entire functions on a
set with an accumulation point; the analytic continuation principle gives
\(F=\gamma_j G\) on all of \(\C\) for any single \(j\). In particular all
\(\gamma_j\) coincide, and \(F\) and \(G\) are linearly dependent, contradicting
our assumption. This proves simplicity.

\medskip
\noindent
{\bf Step 3: a first-variation identity.}

By Step 2, the eigenfunction \(F\) is determined up to a unimodular constant.
We now derive a first-variation identity satisfied by \(F\), based on the fact
that \(\Sigma\) is a level set of \(u_F\) and that holomorphic vector fields are
divergence-free.

The function \(u_F=|F|^2e^{-\pi|z|^2}\) is real-analytic on
\(\C\setminus\{F=0\}\), since \(F\) is entire and the exponential factor is
real-analytic. By Step 2, \(F\) does not vanish in the interior of \(\Sigma\),
and by the argument of Step 4 below it does not vanish on \(\overline\Sigma\)
either; so \(u_F\) is in fact real-analytic on a neighborhood of
\(\overline\Sigma\). The boundary \(\partial\Sigma\) is therefore a level set
of a non-constant real-analytic function and so is the union of finitely many
real-analytic arcs, plus a possible finite set of singular points where
\(\nabla u_F\) vanishes (see, e.g., the standard local structure of
real-analytic varieties). In particular, \(\partial\Sigma\) is Lipschitz, and
\(C^1\) away from the critical set of \(u_F\); the divergence theorem applies
to \(\Sigma\) for any \(C^1\)-vector field on a neighborhood of
\(\overline\Sigma\), the singular boundary points forming a set of
one-dimensional Hausdorff measure zero.

Let \(\psi\) be holomorphic on an open neighborhood \(V\) of
\(\overline \Sigma\), and write
\[
\psi=u+iv
\]
with \(u,v\) real-valued. Consider the real vector field
\[
X_\psi:=(u,-v).
\]
Since \(\psi\) is holomorphic, the Cauchy--Riemann equations
\[
u_x=v_y,
\qquad
u_y=-v_x,
\]
imply
\[
\operatorname{div}X_\psi=u_x+(-v)_y=u_x-v_y=0
\qquad \text{in }V.
\]
Because \(u_F\equiv c_\Sigma\) on \(\partial\Sigma\), the divergence theorem on
\(\Sigma\) gives
\[
\int_\Sigma \langle \nabla u_F, X_\psi\rangle\,dA
=
\int_\Sigma \operatorname{div}(u_F X_\psi)\,dA
=
\int_{\partial \Sigma} u_F\, X_\psi\cdot \nu\, d\mathcal H^1
=
c_\Sigma\int_{\partial \Sigma} X_\psi\cdot \nu\, d\mathcal H^1
=
0,
\]
where the last equality uses
\(\int_{\partial\Sigma}X_\psi\cdot\nu\,d\mathcal H^1
=\int_\Sigma\operatorname{div}X_\psi\,dA=0\) by Cauchy--Riemann. Since
\(u_F=|F|^2e^{-\pi |z|^2}\) and
\(\nabla(e^{-\pi|z|^2})=-2\pi e^{-\pi|z|^2}(x,y)\), this becomes
\begin{equation}\label{eq:vector-field-identity-real-new}
0=
\int_\Sigma
\left(
\langle \nabla |F|^2,X_\psi\rangle
-
2\pi |F|^2 \langle z,X_\psi\rangle
\right)
\,d\mu.
\end{equation}

We now express \eqref{eq:vector-field-identity-real-new} in complex notation.
A direct computation, using \(\partial_z|F|^2=F'\overline F\) and
\(\partial_{\bar z}|F|^2=F\overline{F'}\), gives
\[
\partial_x |F|^2 = F'\overline F+\overline{F'}F = 2\Re(F'\overline F),
\qquad
\partial_y |F|^2 = i(F'\overline F-\overline{F'}F) = -2\Im(F'\overline F).
\]
Writing \(\psi=u+iv\) and recalling \(X_\psi=(u,-v)\),
\[
\langle \nabla |F|^2,X_\psi\rangle
=
2u\,\Re(F'\overline F)+2v\,\Im(F'\overline F)
=
2\,\Re\!\big(\overline{\psi}\,F'\overline F\big)
=
2\,\Re\!\big(\psi\,\overline{F'}F\big),
\]
where in the last equality we used \(\Re w=\Re\overline w\). Likewise, writing
\(z=x+iy\),
\[
\langle z,X_\psi\rangle
=
xu+y(-v)
=
xu-yv
=
\Re\big((x+iy)(u+iv)\big)
=
\Re(z\psi).
\]
Substituting these expressions into
\eqref{eq:vector-field-identity-real-new}, we obtain
\[
0=
\Re\!\int_\Sigma
\psi(z)\Big(\overline{F'(z)}\,F(z)-\pi z\,|F(z)|^2\Big)\,d\mu(z).
\]
Replacing \(\psi\) by \(i\psi\) (which is also holomorphic on \(V\)) gives the
analogous identity for the imaginary part:
\[
0=
\Im\!\int_\Sigma
\psi(z)\Big(\overline{F'(z)}\,F(z)-\pi z\,|F(z)|^2\Big)\,d\mu(z).
\]
Combining the two,
\begin{equation}\label{eq:orthogonality-main-new}
0=
\int_\Sigma
\psi(z)\,F(z)\Big(\overline{F'(z)}-\pi z\,\overline{F(z)}\Big)\,d\mu(z),
\end{equation}
which holds for every holomorphic \(\psi\) on a neighborhood of
\(\overline \Sigma\). Equivalently, after factoring out \(F(z)\) inside the
integrand,
\begin{equation*}
0=
\int_\Sigma
\psi(z)\Big(\overline{F'(z)}\,F(z)-\pi z\,|F(z)|^2\Big)\,d\mu(z).
\end{equation*}
The form \eqref{eq:orthogonality-main-new} is the one used in Step 4.

\medskip
\noindent
{\bf Step 4: \(F'\) is again a first eigenfunction.}

This step is the crux of the argument and is the main new ingredient compared
to classical Faber--Krahn rigidity proofs. Rather than working with general
test functions in the eigenvalue equation, we exploit the divergence identity
\eqref{eq:orthogonality-main-new} together with a clever choice of test
function involving the reproducing kernel of the Fock space. The end result is
that \emph{the derivative} \(F'\) of any first eigenfunction is itself a first
eigenfunction; since the eigenspace is one-dimensional by Step 2, this forces
a first-order ODE on \(F\), and hence pins down \(F\) up to translation as a
Gaussian.

For every \(w\in \C\), set
\[
K_w(z):=e^{\pi \overline w z};
\]
this is the reproducing kernel of \(\mathcal F^2(\C)\) at the point \(w\), in
the sense that
\begin{equation}\label{eq:reproducing-kernel}
\int_\C H(z)\,e^{\pi w \overline z}\,d\mu(z)=H(w),
\qquad H\in\mathcal F^2(\C);
\end{equation}
see, e.g., \cite[Ch.~1]{Folland} or \cite{Bargmann}. Equivalently,
\(\langle H,K_w\rangle_{\mathcal F^2}=H(w)\).

By Step 1, after replacing \(\Sigma\) by the corresponding superlevel-set
representative if necessary, we may suppose that
\[
\Sigma=\{u_F>c_\Sigma\}.
\]
In particular, \(u_F\ge c_\Sigma>0\) on \(\overline \Sigma\). Since
\(u_F(z)=|F(z)|^2e^{-\pi |z|^2}\), it follows that \(F\) has no zeros on
\(\overline \Sigma\). By continuity of \(F\) and the compactness of
\(\overline \Sigma\), there exists an open neighborhood
\(V\supset \overline \Sigma\) on which \(F\) has no zeros. Hence, for every
\(w\in \C\), the function
\[
\psi_w(z):=\frac{K_w(z)}{F(z)}=\frac{e^{\pi\overline wz}}{F(z)}
\]
is holomorphic on \(V\). Applying \eqref{eq:orthogonality-main-new} with
\(\psi=\psi_w\), so that
\(\psi_w(z)F(z)=e^{\pi\overline w z}\), we obtain
\[
0=
\int_\Sigma
e^{\pi \overline wz}
\Big(\overline{F'(z)}-\pi z\,\overline{F(z)}\Big)\,d\mu(z),
\]
or, separating the two terms,
\begin{equation}\label{eq:key-kernel-identity-new}
\int_\Sigma
e^{\pi \overline wz}\overline{F'(z)}\,d\mu(z)
=
\int_\Sigma
\pi z\,e^{\pi \overline wz}\overline{F(z)}\,d\mu(z).
\end{equation}

We now relate the right-hand side of \eqref{eq:key-kernel-identity-new} to
\(\overline{F'(w)}\) using the eigenvalue equation
\eqref{eq:eigenvalue-equation} together with the reproducing identity
\eqref{eq:reproducing-kernel}. Choosing
\[
H(z):=z\,e^{\pi\overline w z}\in \mathcal F^2(\C)
\]
(the product of the kernel \(K_w\) by the polynomial \(z\); both factors
preserve membership in \(\mathcal F^2\)) in \eqref{eq:eigenvalue-equation}, we
obtain
\begin{equation}\label{eq:eigvalue-applied}
\int_\Sigma z\, e^{\pi \overline wz}\,\overline{F(z)}\,d\mu(z)
=
\lambda_{1,\varphi}(\Sigma)\int_\C z\, e^{\pi \overline wz}\,\overline{F(z)}\,d\mu(z).
\end{equation}
On the other hand, taking complex conjugates in \eqref{eq:reproducing-kernel}
applied to \(F\) at the point \(w\), one finds
\[
\int_\C e^{\pi \overline wz}\,\overline{F(z)}\,d\mu(z)=\overline{F(w)}.
\]
Both sides depend on \(w\) through the antiholomorphic variable
\(\overline w\); we differentiate with respect to \(\overline w\), exchanging
differentiation and integration on the left. The interchange is legitimate by
the dominated convergence theorem: locally in \(w\), the integrand
\(\partial_{\overline w}(e^{\pi\overline wz}\overline{F(z)})
=\pi z\,e^{\pi\overline wz}\overline{F(z)}\) is dominated by an
\(L^1(d\mu)\) function on \(\C\), as a consequence of the Cauchy--Schwarz
inequality and the fact that \(F\in \mathcal F^2(\C)\) and the Fock-space norm
of \(zK_w\) is locally bounded. We conclude
\[
\pi\int_\C z\, e^{\pi \overline wz}\,\overline{F(z)}\,d\mu(z)
=
\overline{F'(w)}.
\]
Multiplying \eqref{eq:eigvalue-applied} by \(\pi\) and substituting,
\[
\pi\int_\Sigma z\, e^{\pi \overline wz}\,\overline{F(z)}\,d\mu(z)
=
\lambda_{1,\varphi}(\Sigma)\,\overline{F'(w)}.
\]
Combining this with \eqref{eq:key-kernel-identity-new},
\[
\int_\Sigma e^{\pi \overline wz}\,\overline{F'(z)}\,d\mu(z)
=
\lambda_{1,\varphi}(\Sigma)\,\overline{F'(w)},
\qquad w\in \C.
\]
Taking complex conjugates, this is exactly the eigenvalue equation
\eqref{eq:eigenvalue-equation} applied with \(F\) replaced by \(F'\) (and an
arbitrary \(K_w\) playing the role of the test function). Since the linear
span of \(\{K_w\}_{w\in\C}\) is dense in \(\mathcal F^2(\C)\), the identity
extends to every \(H\in\mathcal F^2(\C)\), and we conclude that \(F'\) is a
first eigenfunction of the localization operator on \(\Sigma\), with eigenvalue
\(\lambda_{1,\varphi}(\Sigma)\).

By the simplicity established in Step 2,
\[
F'=\theta F
\]
for some constant \(\theta\in \C\). Solving this ODE gives
\[
F(z)=C e^{\theta z},
\]
for some nonzero constant \(C\).

\medskip
\noindent
{\bf Step 5: reduction to the constant eigenfunction.}

We now use the Weyl invariance of the Bargmann--Fock setting to reduce the
analysis to the case \(F\equiv 1\). Recall from \eqref{eq:Ta} that, for
\(a\in\C\), the Weyl translation
\[
(T_a F)(z)= e^{\pi a z-\frac{\pi}{2}|a|^2}\,F(z-a)
\]
is unitary on \(\mathcal F^2(\C)\) and satisfies
\(u_{T_aF}(z)=u_F(z-a)\). In particular, \(T_a\) maps eigenfunctions of the
localization operator on \(\Sigma\) to eigenfunctions of the localization
operator on the translated set \(\Sigma+a\), with the same eigenvalue, and the
quantity \(\int_\Sigma u_F\,dA\) is invariant under the simultaneous
substitutions \(F\mapsto T_aF\) and \(\Sigma\mapsto \Sigma+a\).

We claim that we may choose \(a\in\C\) so that \(T_{-a}F\) is a constant
function. Indeed, by Step 4, \(F(z)=Ce^{\theta z}\) for some \(C,\theta\in\C\)
with \(C\ne 0\). Compute, using the identity
\(u_{T_{-a}F}(z)=u_F(z+a)\):
\[
u_{T_{-a}F}(z)=|F(z+a)|^2 e^{-\pi|z+a|^2}
=|C|^2 e^{2\Re(\theta(z+a))}\,e^{-\pi|z+a|^2}.
\]
Expanding,
\[
2\Re(\theta(z+a))-\pi|z+a|^2
=
2\Re(\theta z)+2\Re(\theta a)-\pi|z|^2-2\pi\Re(\overline a z)-\pi|a|^2.
\]
The linear terms in \(z\) on the right cancel if and only if
\(2\Re(\theta z)=2\pi\Re(\overline a z)\) for every \(z\in\C\), which is the
condition \(\overline a=\theta/\pi\), i.e.
\[
a:=\frac{\overline\theta}{\pi}.
\]
With this choice,
\[
u_{T_{-a}F}(z)=|C'|^2 e^{-\pi|z|^2},
\qquad
|C'|^2:=|C|^2 e^{2\Re(\theta a)-\pi|a|^2}=|C|^2 e^{|\theta|^2/\pi}.
\]
Since \(T_{-a}F\) is entire and \(|T_{-a}F(z)|^2 e^{-\pi|z|^2}=|C'|^2
e^{-\pi|z|^2}\), it follows that \(|T_{-a}F|\equiv|C'|\) is constant; a
holomorphic function with constant modulus is constant, hence
\(T_{-a}F\equiv C''\) for some \(C''\in\C\) with \(|C''|=|C'|\). Normalizing,
we may absorb the unimodular constant into the eigenfunction and assume
\[
F\equiv 1
\]
after the change \(\Sigma\rightsquigarrow \Sigma-a\) (and the corresponding
unitary substitution \(F\rightsquigarrow T_{-a}F\)).

In this normalization the ground state \(h_0\) (corresponding to
\(F\equiv 1\) in Bargmann--Fock coordinates) is an eigenfunction of the
localization operator of \(\Sigma\). Since \(\Sigma\) is a level set of
\(u_F(z) = e^{-\pi |z|^2}\), and the latter is a strictly radial function of
\(|z|\), the set \(\Sigma\) is itself radial; in particular, it is either a
centered disk or a finite union of centered annuli. The next step rules out
the annular configurations and identifies \(\Sigma\) with the disk.

\medskip
\noindent
{\bf Step 6: exclusion of annuli and conclusion.}

It remains to verify that, among radial sets of prescribed measure, the only
local maximizer is the centered disk. Although this is essentially automatic
after Step 5, we record the standard rearrangement argument since it makes
explicit \emph{why} no annular local maximizer can exist; this provides the
endpoint matching the global Faber--Krahn theorem of \cite{Nicola-Tilli}.

If \(\Sigma\) is radial, write
\[
\widetilde \Sigma:=\{r>0:\ re^{it}\in \Sigma \text{ for some } t\in [0,2\pi)\}.
\]
By Step 1 and the Lipschitz regularity established in Step 3,
\(\widetilde \Sigma\) is open and is a finite union of disjoint open intervals:
\[
\widetilde \Sigma=\bigcup_{j=1}^N (a_j,b_j),
\qquad
0\le a_1<b_1\le a_2<b_2\le \cdots\le a_N<b_N.
\]
Then, using \(\int re^{-\pi r^2}\,dr=-\frac{1}{2\pi}e^{-\pi r^2}\),
\[
\int_\Sigma e^{-\pi |z|^2}\,dA(z)
=
2\pi \int_{\widetilde \Sigma} r e^{-\pi r^2}\,dr
=
\sum_{j=1}^N \big(e^{-\pi a_j^2}-e^{-\pi b_j^2}\big),
\]
under the constraint
\[
|\Sigma|
=
2\pi \int_{\widetilde \Sigma} r\,dr
=
\pi \sum_{j=1}^N (b_j^2-a_j^2).
\]
We are therefore led to the optimization problem
\begin{equation}\label{eq:radial-opt}
\text{maximize}\qquad
\sum_{j=1}^N \big(e^{-\pi a_j^2}-e^{-\pi b_j^2}\big)
\qquad \text{subject to}\qquad
\sum_{j=1}^N (b_j^2-a_j^2)=c,
\end{equation}
where \(c=|\Sigma|/\pi\) is fixed. Substituting \(s_j:=a_j^2\),
\(t_j:=b_j^2\) (with \(0\le s_1<t_1\le s_2<\cdots\)), and writing
\(\Phi(x):=e^{-\pi x}\), \eqref{eq:radial-opt} becomes
\[
\text{maximize}\qquad
\sum_{j=1}^N \big(\Phi(s_j)-\Phi(t_j)\big)
\qquad \text{subject to}\qquad
\sum_{j=1}^N (t_j-s_j)=c.
\]
The function \(\Phi\) is strictly convex and strictly decreasing on
\([0,\infty)\); its derivative \(-\Phi'(\sigma)=\pi e^{-\pi\sigma}\) is strictly
decreasing in \(\sigma\). For any \(t>s\ge 0\),
\[
\Phi(s)-\Phi(t)
=
\int_s^t (-\Phi'(\sigma))\,d\sigma,
\]
so
\[
\sum_{j=1}^N\big(\Phi(s_j)-\Phi(t_j)\big)
=
\int_{\widetilde\Sigma_2}(-\Phi'(\sigma))\,d\sigma,
\qquad
\widetilde\Sigma_2:=\bigcup_{j=1}^N (s_j,t_j).
\]
The set \(\widetilde\Sigma_2\subset[0,\infty)\) has Lebesgue measure equal to
\(c\), and the integrand \(-\Phi'\) is strictly decreasing in \(\sigma\). The
bathtub principle (or, equivalently, a direct rearrangement argument) shows
that, among all measurable subsets of \([0,\infty)\) of total Lebesgue measure
\(c\), the integral \(\int_{\widetilde\Sigma_2}(-\Phi'(\sigma))\,d\sigma\) is
strictly maximized by the leftmost interval \((0,c)\). This corresponds to
\(N=1\), \(s_1=0\), \(t_1=c\), i.e.
\[
\widetilde \Sigma=(0,R),
\qquad R:=\sqrt c,
\]
and equivalently \(\Sigma\) is the centered disk \(D_R\) with
\(|D_R|=|\Sigma|\). Any annular configuration (\(N\ge 2\), or \(N=1\) with
\(s_1>0\)) produces a strictly smaller value of \(\int_\Sigma u_F\,dA\), and is
therefore not a local maximizer.

We conclude that every local maximizer of \(\lambda_{1,\varphi}\) under the
measure constraint is, up to translation, a disk.
\end{proof}

